% Examen final -- EPD-GIT-5534 (6 exercices, 1 h 30)
\documentclass[11pt,a4paper]{article}
\usepackage[margin=1.8cm]{geometry}
\usepackage[french]{babel}
\usepackage{amsmath,amssymb}
\usepackage{graphicx}
\usepackage{booktabs}
\usepackage{enumitem}
\usepackage{xcolor}
\usepackage{tcolorbox}
\usepackage{fancyhdr}
\usepackage{lastpage}
\usepackage[hidelinks]{hyperref}
\newcommand{\py}[1]{\texttt{\small #1}}
\newcommand{\file}[1]{\texttt{\small #1}}
\definecolor{bleuenspd}{HTML}{1F4E79}
\pagestyle{fancy}\fancyhf{}
\fancyhead[L]{\small EPD-GIT-5534 --- Apprentissage automatique appliqué aux infrastructures}
\fancyhead[R]{\small\thepage/\pageref{LastPage}}
\fancyfoot[C]{\small Durée : 1 h 30 --- documents non autorisés --- calculatrice autorisée}

\newtcolorbox{donnees}{colback=gray!6, colframe=gray!55, boxrule=0.5pt, arc=1.5mm,
  left=2mm, right=2mm, top=1.5mm, bottom=1.5mm, title={Rappels et données},
  fonttitle=\bfseries\small, coltitle=white, colbacktitle=gray!70}

\setlist{nosep, leftmargin=*}

\begin{document}

\begin{center}
\includegraphics[width=2.4cm]{../assets/logo_enspd.png}\\[1mm]
{\large\color{bleuenspd}\textbf{École Nationale Supérieure Polytechnique de Douala}}\\
Université de Douala --- Département Génie Informatique et Télécommunications\\[2mm]
{\Large\color{bleuenspd}\textbf{Examen final --- Apprentissage automatique appliqué aux infrastructures}}\\[1mm]
EPD-GIT-5534 --- Master Professionnel ASIC --- Durée : \textbf{1 h 30} --- \textbf{20 points}\\
Enseignante : Dr. Noulapeu Ngaffo Armielle --- Session : Normale
\end{center}

\begin{donnees}
Le dossier NOSYS tout au long de l'épreuve : 8 serveurs (\py{srv-01} à \py{srv-08}),
30 jours de métriques à pas de 5 min (69 120 fenêtres), SLO « latence p99 $<$ 100 ms »,
SLA de disponibilité 99,9 \%. Chiffres mesurés sur le mois :
\begin{itemize}
\item CPU moyen du parc 20,5 \% (heures ouvrées 30,2 \%, nuit 10,0 \%) ; CPU max incident 96 \% ;
\item latence : moyenne 12,2 ms, p99 16,3 ms ; température max 42,7 °C ; consommation 11 614 kWh ;
\item 16 incidents (2 par serveur : 35 min le 09/09 à 14 h, 50 min le 22/09 à 3 h),
soit 11,3 h d'indisponibilité cumulée sur $8 \times 720$ h ;
\item dataset supervisé : 34 560 fenêtres de 10 min, 120 positives (0,35 \%) ;
train (21 j) = 24 192 lignes dont 56 positives ; test (9 j) = 10 368 lignes dont 64 positives ;
\item logs : 100 000 lignes d'authentification sur 7 jours, 4 828 échecs,
dont une attaque brute force le 05/10 de 2 h à 4 h (125 échecs sur \py{bdd-principal}).
\end{itemize}
Formules : $z = (x-\bar{x})/\sigma$ ; IQR $= Q_3 - Q_1$ ; précision $= VP/(VP+FP)$ ;
rappel $= VP/(VP+FN)$ ; $F_1 = 2\,\text{précision}\cdot\text{rappel}/(\text{précision}+\text{rappel})$ ;
disponibilité $= 1 - \text{indispo}/\text{temps total}$.
\end{donnees}

\section*{Exercice 1 --- Cadre et indicateurs (3 pts)}
a) (0,5) Donner la différence entre SLI, SLO et SLA, avec un exemple de chacun pour NOSYS.\\
b) (0,5) Calculer la disponibilité du mois de NOSYS à partir des 11,3 h d'indisponibilité
(8 serveurs $\times$ 720 h). Conclure par rapport au SLA.\\
c) (1) Le p99 de latence est de 16,3 ms. Expliquer pourquoi on supervise le p99 plutôt
que la moyenne, en quantifiant l'impact sur 10 millions de requêtes/jour.\\
d) (1) Classer ces 4 problèmes dans la bonne famille de ML (supervisé régression /
supervisé classification / non supervisé / renforcement) :
(i) prévoir le trafic à 24 h ; (ii) détecter des sessions d'auth anormales sans
étiquettes ; (iii) prédire un incident dans les 30 min ; (iv) piloter le refroidissement
par essai-récompense.

\section*{Exercice 2 --- Préparation des données (3 pts)}
Le fichier d'acquisition brut contient : 200 lignes dupliquées, 5 508 valeurs manquantes
réparties uniformément (1,73 \%), une température de 1 575 °C, un CPU à 350 \%,
6 trous de 15 à 40 min, et un serveur qui relève toutes les 15 min.
a) (0,75) Classer le mécanisme des manquants (MCAR/MAR/MNAR) et justifier en une phrase.\\
b) (0,75) CPU à 350 \% : erreur de mesure ou incident ? Même question pour un CPU à 96 \%
pendant une alarme. Quelle règle de traitement en déduire ?\\
c) (0,75) Pourquoi trier par (serveur, horodatage) avant toute moyenne mobile ?\\
d) (0,75) Donner deux raisons d'encoder l'heure avec (sin, cos) plutôt qu'une colonne 0--23.

\section*{Exercice 3 --- Détection d'anomalies et classe rare (3,5 pts)}
Le dataset supervisé compte 34 560 fenêtres dont 120 positives (0,35 \%).
Un modèle prédit « incident » sur 60 fenêtres du test ; parmi elles, 40 sont de vrais
incidents ; il manque 24 vrais incidents.
a) (1) Compléter la matrice de confusion (VP, FP, FN, VN avec VN $= 10\,368 - 60$)
et calculer précision, rappel et F1.\\
b) (0,75) Un modèle « toujours 0 » atteindrait quelle exactitude sur le test ? Pourquoi
cette métrique est-elle trompeuse ici ?\\
c) (1) Pourquoi la courbe précision-rappel est-elle préférée à la ROC sur ce problème ?
Que signifie un point d'opération à (précision 1,0 ; rappel 0,625) pour le NOC ?\\
d) (0,75) Citer deux raisons pour lesquelles un seuil fixe « CPU $>$ 90 \% » est
fragile, même s'il a détecté 136/136 fenêtres ce mois-ci.

\section*{Exercice 4 --- Série temporelle et prédiction (3,5 pts)}
On agrège la charge horaire (720 points). Le test ADF donne $p = 0{,}075$ sur la série
brute et $p = 3{,}4\cdot10^{-7}$ après différenciation d'ordre 24.
Un SARIMA $(2,0,1)\times(1,0,1)_{24}$ donne une MAE de 3,01 points de CPU sur 7 jours,
contre 5,06 pour la persistance et 5,23 pour « la même semaine la semaine dernière ».
a) (1) Interpréter les deux valeurs de l'ADF. Quel rôle joue la différenciation d'ordre 24 ?\\
b) (0,75) Calculer le gain du SARIMA par rapport à la persistance en \%.\\
c) (1) Expliquer pourquoi une validation croisée KFold mélangée est invalide ici, et
décrire le protocole correct (train/test par dates).\\
d) (0,75) La fuite de données : un étudiant ajoute la colonne \py{alarme} courante aux
variables. Mesurée : P(incident $\mid$ alarme = 1) = 100 \%. Expliquer le résultat obtenu
en validation et pourquoi il s'effondre en production.

\section*{Exercice 5 --- Sécurité (3,5 pts)}
Sur 7 jours de logs (100 000 lignes), on agrège par (machine, fenêtre de 10 min) :
10 000 sessions, dont 12 sessions d'attaque (brute force du 05/10, 2 h--4 h,
125 échecs, 0 succès, uniquement comptes \py{root}/\py{admin}/\py{inconnu-*}).
Une IsolationForest classe les sessions par score d'anomalie ; les 12 sessions
d'attaque occupent les rangs 1 à 12.
a) (1) Calculer la précision@10 et la précision@30 de l'isolation forest (les sessions
d'attaque occupent les rangs 1--12). Que vaut la précision attendue au hasard à rang 10 ?\\
b) (1) Proposer 4 variables (features) d'authentification utiles pour détecter cette
attaque, en justifiant chacune en une ligne.\\
c) (0,75) L'entropie des comptes testés vaut 1,73 bits pendant l'attaque contre 2,22 bits
en normal. Interpréter.\\
d) (0,75) Citer deux limites d'un détecteur ML face à un attaquant adaptatif, et une parade.

\section*{Exercice 6 --- Décision, capacité et industrialisation (3,5 pts)}
a) (1,5) Un modèle prédictif permet d'éviter 70 \% des incidents (détectés 30 min avant).
Recalculer la disponibilité du mois (indisponibilité initiale 11,3 h) et conclure par
rapport au SLA de 99,9 \%.\\
b) (1) Le stockage est à 60 \% ; la consommation de données croît de 1,333 \%/jour.
Calculer la date de saturation (seuil d'alerte 90 \%), en indiquant la formule.\\
c) (1) Le modèle de production score une fenêtre en 59 ms en moyenne (p99 : 105 ms).
Combien de fenêtres peut scorer un seul cœur par minute ? Quel monitoring mettre en place
(3 métriques minimales) pour détecter une dérive du modèle en production ?

\vspace{2mm}
\begin{center}\small
\textbf{Barème :} Ex. 1 : 3 pts --- Ex. 2 : 3 pts --- Ex. 3 : 3,5 pts --- Ex. 4 : 3,5 pts ---
Ex. 5 : 3,5 pts --- Ex. 6 : 3,5 pts. \quad Total : 20 pts.\\
Tous les chiffres nécessaires figurent dans l'encadré « Rappels et données » ou l'énoncé ;
tout calcul doit être posé.
\end{center}

\end{document}
